\subsection{Integral extensions from pointed Morita data}
\label{sec:extension-integrality}
\begin{lemma}\label{lem:square-kernels}
Suppose $\mathcal B$ is Morita equivalent to $\Vec_G^\omega$ and every
normal abelian subgroup of $G$ has square order. Then every faithful
finite-group graded extension of $\mathcal B$ is integral.
\end{lemma}
\begin{proof}
An invertible $\Vec_G^\omega$ bimodule is, as a left module, a transitive
pointed module associated to a subgroup $N\leq G$ and a cochain
trivializing $\omega|_N$. Its dual is pointed
\cite[Proposition~4.2]{ENOExt}. In the twisted double-coset description
\cite[Theorems~4.1 and 5.1]{GelakiNaidu}, a simple over $NaN$ has
FP dimension
\[
[N:N\cap aNa^{-1}]\dim(\pi),
\]
where $\pi$ is an irreducible projective representation of the
intersection. Pointedness forces $N$ to be normal. The identity
coset subcategory is ordinary $\Rep(N)$, even with the general
associator and subgroup cochain \cite[Theorem~4.1]{GelakiNaidu}; thus
$N$ is abelian as well. The left module has $[G:N]$ simples with equal
normalized FP dimension. Since a graded component has total squared
dimension $|G|$, each dimension is $\sqrt{|N|}$, an integer.
Every component of an extension of $\Vec_G^\omega$ is an invertible
bimodule, so that extension is integral.

Morita equivalence transports the complete graded-extension data,
including multiplication and associativity, and the resulting complete
extensions are Morita equivalent
\cite[Corollary~3.13 and Theorem~1.3]{ENOExt}.
Integrality is preserved under this equivalence in the present setting.
Indeed, the center of an integral category is integral because
forgetting preserves FP dimensions. A Morita equivalent category
$\mathcal C$ has the same center and integer total dimension. For a
simple $W\in\mathcal C$, induction has dimension
$\FPdim(\mathcal C)\FPdim(W)$, an integer. Hence $\FPdim(W)$ is rational
and an algebraic integer, and therefore is an integer.
\end{proof}

\begin{proposition}\label{prop:dimension72}
The rings $R8\_2\_53$ and $R8\_2\_93$ have no categorification over $\C$.
\end{proposition}
\nopagebreak[4]
\begin{proof}
Each source tensor has FP type $1^4,4^2,(3\sqrt2)^2$ and a faithful
$C_2$ grading with integral neutral subring $D$ of dimension $36$.
Its invertibles form $C_4=\langle g\rangle$. They fix both
noninvertibles $X,Y$, which are self-dual, and
\[
X^2=1+g+g^2+g^3+X+2Y,\qquad
Y^2=1+g+g^2+g^3+2X+Y.
\]
In particular $D_{\rm ad}=D$. Any categorification $\mathcal D$ is
solvable by the two-prime theorem \cite[Theorem~1.6]{ENOSolvable}.
Its last prime step must be an equivariantization since its universal
grading is trivial. The prime cannot be three because
$\operatorname{Inv}(\mathcal D)=C_4$ contains no $C_3$. Thus
$\mathcal D=\mathcal E^{C_2}$, with $\FPdim(\mathcal E)=18$.
The sign object is $g^2$ and fixes $X,Y$. The equivariantization
orbit formula \cite{BurciuNatale} makes $\mathcal E$ of type $1^2,2^4$.

Let $\eta$ be its nontrivial invertible. In every $ZZ^*$ with
$\FPdim(Z)=2$, omitting $\eta$ would leave dimension three after the
unit, which cannot be a sum of two-dimensional simples. Therefore
$ZZ^*=1+\eta+T$ with $\FPdim(T)=2$. The adjoint subcategory contains
both units and a two-dimensional simple. If proper, its integer
dimension divides $18$; the only possibilities for a subcategory of
type $1^2,2^n$ are $2$ and $6$. Dimension two omits $T$. At dimension
six, the other universal-grading components contain only
two-dimensional simples and cannot have total squared dimension six.
Hence $\mathcal E_{\rm ad}=\mathcal E$.
Solvability again forces $\mathcal E=\mathcal P^{C_2}$, where
$\FPdim(\mathcal P)=9$ and $\mathcal P$ is pointed \cite{ENO,EGNO}.

The canonical central $\Rep(C_2)\subset Z(\mathcal E)$ has underlying
object $\eta$. Its half-braiding is unique: two half-braidings on an
invertible differ by a tensor character, and
$\operatorname{Aut}_\otimes(\operatorname{Id}_{\mathcal E})
=\operatorname{Hom}(U(\mathcal E),\C^\times)=1$.
The outer $C_2$ action therefore preserves this central subcategory.
Its equivariantization embeds as a symmetric subcategory of
$Z(\mathcal E^{C_2})$: equivariant half-braidings act on every
equivariant object, and the morphism conditions give full faithfulness.
Forgetting equivariance followed by the fiber functor on $\Rep(C_2)$
gives a fiber functor, so this subcategory is $\Rep(H)$ for a group
of order four. Its four underlying invertibles exhaust the $C_4$
in $\mathcal D$, forcing $H=C_4$.
Central de-equivariantization \cite{DGNO,EGNO} now gives
\[
\mathcal D\simeq(\Vec_A^\alpha)^{C_4},\qquad |A|=9.
\]
The dimension-four noninvertibles force $C_4$ to act freely on all
nonidentity elements of $A$. Thus $A=C_3^2$, since
$\operatorname{Aut}(C_9)$ has order six. Its generator acts by an
irreducible matrix $J$ over $\mathbb F_3$ with $J^2=-I$.

The pointed crossed product is Morita dual to this equivariantization.
Consequently $\mathcal D$ is Morita equivalent to $\Vec_G^\omega$,
where $G=C_3^2\rtimes_J C_4$; the group extension splits by coprimality.
The only normal abelian subgroups of $G$ are $1$ and $C_3^2$.
Indeed, an element with nonzero $C_4$ exponent has commutators with
$C_3^2$ filling $(J^k-I)C_3^2=C_3^2$. A normal subgroup containing
such an element contains $C_3^2$ and cannot be abelian. Within
$C_3^2$, irreducibility of $J$ leaves only $1$ and $C_3^2$.
Their orders are squares, so Lemma~\ref{lem:square-kernels} makes
every extension of $\mathcal D$ integral. The alleged census category
would be a $C_2$ extension with two simples of dimension $3\sqrt2$,
a contradiction. The exact verifier checks both full source tensors,
the grading and every finite-group premise.
\end{proof}

\begin{proposition}\label{prop:dimension24}
The rings $R8\_2\_191$ and $R8\_2\_299$ have no categorification over $\C$.
\end{proposition}
\nopagebreak[4]
\begin{proof}
Their source tensors have FP type $1^4,2^2,(\sqrt6)^2$ and a faithful
$C_2$ grading with integral neutral subring $D$ of dimension $12$.
Its invertibles form $V_4$ and fix two simples $X,Y$, with
\[
X^*=Y,\qquad X^2=2Y,\qquad Y^2=2X,\qquad
XY=YX=\sum_{g\in V_4}g.
\]
Thus a hypothetical neutral category $\mathcal D$ is faithfully
$C_3$ graded over its pointed subcategory $\mathcal B$ on $V_4$,
with rank-one nonneutral components. A rank-one module trivializes
the pointed associator, so $\mathcal B\simeq\Vec_{V_4}$.

The ambient odd component $\mathcal M$ has two simples permuted
transitively by $V_4$, with a stabilizer $H=C_2$; its stabilizer in
source labels is $(0,5)$ for $R8\_2\_191$ and $(0,7)$ for
$R8\_2\_299$. Its restriction $\mathcal N$ to $\mathcal B$ is
therefore an indecomposable rank-two module. Since
$H^2(C_2,\C^\times)=0$, its cochain is trivial. A decomposition
$V_4=H\times K$ identifies its dual with
\[
\mathcal B^*_{\mathcal N}
\simeq\Rep(H)\boxtimes\Vec_K\simeq\Vec_{V_4},
\]
including the trivial associator.
The actual multiplication action of $\mathcal D$ on $\mathcal M$
is a coherent extension of this $\mathcal B$ module.

Apply \cite[Theorem~1.3]{GalindoClifford} to the \emph{induced} module
$\mathcal L=\mathcal D\boxtimes_{\mathcal B}\mathcal N$.
It makes the dual $\mathcal D^*_{\mathcal L}$ a semidirect $C_3$
product with neutral component $\mathcal B^*_{\mathcal N}$, and
hence pointed. The theorem concerns this induced-module dual;
it does not assert that the dual of the original two-object
$\mathcal M$ is pointed.

The original $C_3$ Brauer--Picard action on $\mathcal B$ is nontrivial:
a rank-one component is not the regular rank-four module. Its class
therefore has order three. Morita transport conjugates the action and
preserves this order. The tensor-automorphism classes of $\Vec_{V_4}$
have kernel $H^2(V_4,\C^\times)=C_2$ over
$\operatorname{Aut}(V_4)=S_3$ \cite{EGNO}. An order-three class must
therefore act nontrivially on $V_4$. The pointed Morita group is
$V_4\rtimes C_3=A_4$, so $\mathcal D$ is Morita equivalent to
$\Vec_{A_4}^\omega$ for some $\omega$.
The only normal abelian subgroups of $A_4$ are $1$ and $V_4$, of
square orders $1$ and $4$. Lemma~\ref{lem:square-kernels} makes every
extension of this compatible neutral category integral, contradicting
the odd FP dimensions $\sqrt6$. The source checker verifies both
gradings, the complete neutral tensors, the odd module action and
all ten subgroups of $A_4$. The neutral ring by itself is not claimed
to have only pointed Morita models; the coherent odd module is essential.
\end{proof}

\begin{proposition}\label{prop:dimension48}
The rings $R8\_2\_158$, $R8\_2\_186$, $R8\_2\_579$ and
$R8\_2\_682$ have no categorification over $\C$.
\end{proposition}
\nopagebreak[4]
\begin{proof}
All four source tensors have FP type $1^4,2,4,(2\sqrt3)^2$ and a
faithful $C_2$ grading. Their common integral neutral subring $D$,
of dimension $24$, has $V_4$ invertibles fixing $Y,Z$ and
\[
Y^*=Y,\quad Z^*=Z,\qquad Y^2=\sum_{g\in V_4}g,
\quad YZ=ZY=2Z,\quad Z^2=\sum_{g\in V_4}g+2Y+2Z.
\]
Thus $D_{\rm ad}=D$. Solvability \cite[Theorem~1.6]{ENOSolvable}
forces a neutral categorification to be
$\mathcal D=\mathcal E^{C_2}$ with $\FPdim(\mathcal E)=12$;
its invertibles exclude a final $C_3$ equivariantization.
The sign $\epsilon$ fixes $Y,Z$. The orbit formula makes
$\mathcal E$ of type $1^4,2^2$, with an outer involution swapping
the two noninvertibles and swapping a pair of invertibles.

The finite rank-six check exhausts both groups $A=C_4,V_4$, their
trivial and transitive actions on the two nonunits, both dualities,
and every oriented nonunit product allowed by FP dimension and
reciprocity. There are $100$ candidate tables and six associative
ones, of which exactly two admit the required outer involution.
Their rings are $R_{3,A}$:
\[
X^*=X',\qquad X^2=2X',\qquad (X')^2=2X,\qquad
XX'=X'X=\sum_{g\in A}g,
\]
with all units fixing $X,X'$. The transitive $C_4$ action in this
exhaustion is the homomorphism $(0,1,0,1)$ in cyclic labels.
Jordan--Larson \cite[Theorem~1.4]{JordanLarson} excludes $A=C_4$,
since its cyclic prime-power factor occurs with odd multiplicity.
Thus $\mathcal E$ is $C_3$ graded over
$\mathcal B\simeq\Vec_{V_4}$; its rank-one components trivialize
the pointed associator.

Let $\mathcal M$ be the two-object ambient odd module over
$\mathcal D$. The central regular algebra $1+\epsilon$ in
$\mathcal D$ gives an actual de-equivariantized module
$\mathcal M_{1+\epsilon}$ over $\mathcal E$.
This construction uses the central algebra inside $\mathcal D$;
it requires no half-braiding on $\epsilon$ in the ambient category.
Its rank is one if $\epsilon$ swaps the source simples, and four
if it fixes them: each fixed simple has two algebra-module structures,
using $H^2(C_2,\C^\times)=0$.

The six indecomposable $\Vec_{V_4}$ modules have ranks
$4,1,1,2,2,2$, by their stabilizers and cochains.
The $C_3$ action sends the regular rank-four class to the $X$
component, then to the $X'$ component. These are distinct rank-one
classes, so they form a three-cycle of total rank six. The three
rank-two classes are either fixed or form another orbit of total
rank six. Galindo's Clifford theorem
\cite[Proposition~4.1 and Theorem~1.1]{GalindoClifford} now excludes
rank one for $\mathcal M_{1+\epsilon}$, and makes its rank-four
alternative two fixed rank-two submodules. The group $C_3$ cannot
permute two actual summands; each is a coherent $\mathcal E$-module
extension of a rank-two $\mathcal B$ module.

For one such module, the induced-module dual of
\cite[Theorem~1.3]{GalindoClifford} is a semidirect $C_3$ product
with neutral $\Vec_{V_4}$. Its action class has order three, as in
Proposition~\ref{prop:dimension24}.
Since the tensor-automorphism classes are $C_2\times S_3$, this
class is a pure order-three group action. The remaining action
coherence lies in $H^2(C_3,\widehat{V_4})=0$ by coprimality
\cite[Proposition~3.2]{NikshychRiepel}. Hence the action is equivalent
to the strict group action and the dual is the untwisted
$\Vec_{A_4}$. Its graded tensor section is part of Galindo's
semidirect-product conclusion; no uncontrolled $C_3$ associator
twist is discarded.

The equivariantization $\mathcal D$ is Morita equivalent to
$\mathcal E\rtimes C_2$. Morita transport gives a $C_2$ extension
$\mathcal T$ of $\Vec_{A_4}$ with the same center as $\mathcal D$.
The center of $\Vec_{A_4}$ has exactly three Lagrangian subcategories,
all in the orbit of the canonical $\Rep(A_4)$
\cite[Proposition~7.6 and Section~8.3]{NikshychRiepel}.
Consequently each has untwisted $\Vec_{A_4}$ as pointed Morita dual.
The involution fixes at least one of these three. For its invariant
pointed module $\mathcal N$, the dual of
$\mathcal T\boxtimes_{\Vec_{A_4}}\mathcal N$ is a $C_2$ crossed
product with pointed neutral $\Vec_{A_4}$
\cite[Proposition~5.3]{GalindoClifford}; every component is pointed.
This step requires module-class invariance, rather than a coherent
extension of $\mathcal N$. It gives a pointed Morita group
\[
1\longrightarrow A_4\longrightarrow G\longrightarrow C_2
\longrightarrow1.
\]
Because $A_4$ is centerless with $\operatorname{Out}(A_4)=C_2$,
$G$ is $S_4$ or $A_4\times C_2$.
The latter has six identity-supported central characters for every
associator. But
$\operatorname{Inv}(Z(\mathcal D))\to\operatorname{Inv}(\mathcal D)$
is injective: its kernel is
$\operatorname{Hom}(U(\mathcal D),\C^\times)=1$.
The four source invertibles therefore exclude $A_4\times C_2$.
Thus $\mathcal D$ is Morita equivalent to $\Vec_{S_4}^\omega$.

The only normal abelian subgroups of $S_4$ are $1$ and its $V_4$.
Lemma~\ref{lem:square-kernels} makes every extension of
$\mathcal D$ integral, contradicting the two odd dimensions
$2\sqrt3$. The exact checker validates all four source tensors,
the small-ring exhaustion, module-rank alternatives and all thirty
subgroups of $S_4$.
\end{proof}
